\section{Problem Formulation}
\label{sec:problem}

\paragraph{Setting.}
At decision time $t$, a frozen image encoder $f$ maps the robot's camera images to patch features.
The context $C$ collects the features of the current and the previous image together with recent proprioception; it is computed once per decision.
A candidate action chunk $\bA=(a_t,\dots,a_{t+\He-1})$, if executed, would produce a \emph{future segment} $\tau$: the features of the images it generates, subsampled in time, together with the final proprioception.
The world model's prediction about $\tau$ must suffice to rank the candidates before any of them is executed.

\paragraph{Selecting among policy proposals.}
At each decision, a frozen pretrained policy $\pi$ proposes $K$ chunks, where $\bA^{(1)}$ is the chunk the policy would execute on its own.
A scorer ranks the proposals, the highest-ranked chunk is executed, with ties resolved in favor of $\bA^{(1)}$, and the policy replans after $\He$ steps~\cite{qi2025gpc,cui2026dreamsteer}.
Optimizing actions against a learned model tends to exploit its errors~\cite{gao2023overoptimization}; selecting among the proposals of a competent policy instead keeps the candidates close to the data on which the model was trained~\cite{janner2019mbpo}.
All scorers in this paper rank the same proposals.

\paragraph{Goals and labels.}
Goals are specified by images $g\in\cG$.
For training, the simulator labels each candidate with the task progress $\ell$ reached at the end of its segment; on PushT, $\ell$ is the negative mean distance between the block's vertices and their target positions, which orders candidates in the same way as the vertex-registration reward of GPC~\cite{qi2025gpc}.
Candidates of one decision share their starting state, so only their order matters.
No label is used at test time.

\paragraph{Conditional trajectory code.}
A target encoder $q_\phi$, used only during training, maps the context and the actual future to $M$ discrete tokens, $S=q_\phi(C,\tau)\in\{1,\dots,V\}^M$.
Two heads read the code together with the context, but never the action: a reader $D_\psi(C,S,g)$ that scores progress toward $g$, and a feature decoder $G_\xi(C,S)$ that reconstructs $\tau$.
Given $K$ sibling futures $\tau_1,\dots,\tau_K$ of one decision with labels $\ell_1,\dots,\ell_K$ and codes $S_k=q_\phi(C,\tau_k)$, the code solves the conditional compression problem
\begin{equation}
\begin{aligned}
  \min_{\phi,\psi,\xi}\ &\E\Big[\mathcal{L}_{\mathrm{rank}}\big(\{D_\psi(C,S_k,g)\}_k,\{\ell_k\}_k\big)\\
  &\qquad+\frac{\alpha}{K}\sum\nolimits_{k} d_{\mathrm{feat}}\big(G_\xi(C,S_k),\tau_k\big)\Big]\\
  \text{s.t.}\ &I(\tau;S\mid C)\le R,
\end{aligned}
\label{eq:crd}
\end{equation}
where the expectation is over decisions and goals, $\mathcal{L}_{\mathrm{rank}}$ is a ranking loss over the siblings (\cref{sec:reader}), and $d_{\mathrm{feat}}$ is a reconstruction error.
Both the encoder and the heads see the context, so the code need not repeat what $C$ already contains.
Since $S$ is a deterministic function of $(C,\tau)$, $I(\tau;S\mid C)=H(S\mid C)\le M\log_2 V$; we set the rate through the code length, and $M{=}16$ tokens with $V{=}256$ values bound it at 128 bits per segment.

A world model $p_\theta(S\mid C,\bA)$ must forecast the code before execution.
Given the context, the action can reduce the uncertainty about the code by at most $I(S;\bA\mid C)$; the remainder, $H(S\mid C,\bA)$, reflects dynamics that the context and the action do not determine and cannot be forecast.
We estimate $I(S;\bA\mid C)$ in \cref{sec:wm}.

\paragraph{Two limiting cases.}
The family in \eqref{eq:crd} contains two familiar designs as extremes (\cref{tab:designs}).
A \emph{per-step latent world model} such as DINO-WM sets $S=\tau$: it predicts every token of every future frame, including everything the history already determines, and ranks with a fixed feature distance or a learned readout.
A \emph{direct scorer} $D(C,\bA,g)$ sets $S=\bA$: nothing about the future is predicted, and the reader must learn the dynamics together with the evaluation.
\method places $S$ between these extremes, as a learned description of the future at a bounded rate that the reader can evaluate for any goal.

\begin{table}[t]
\centering
\small
\setlength{\tabcolsep}{2pt}
\begin{tabular}{@{}llcc@{}}
\toprule
 & Predicted & Sequential & Supervised \\
 & per candidate & passes & by future \\
\midrule
Per-step latent WM & $\He P$ feature tokens & $\He$ & \checkmark \\
Endpoint latent WM & $P$ feature tokens & 1 & \checkmark \\
Direct scorer & nothing ($S=\bA$) & 1 & -- \\
\method (ours) & $M$ code tokens & 1 & \checkmark \\
\bottomrule
\end{tabular}
\caption{\textbf{What each design predicts} for one candidate chunk, with $P{=}256$ patch tokens per image, $M{=}16$ code tokens, and $\He{=}8$ on PushT. Except for the direct scorer, the prediction does not depend on the goal and is reused for every goal; only the readout is repeated.}
\label{tab:designs}
\end{table}
